\subsection{Auxiliary definitions}

The operators and the other auxiliary functions the evaluation rules use are defined in
\secref{auxiliary-values}, and module and class resolution in \secref{auxiliary-classes}.

\begin{AddedBlock}
\subsubsection{Operators}\label{sec:run-time-overloads}

\begin{figure}
\small
\centering
\begin{tabular}{>{\raggedright\arraybackslash}p{0.14\textwidth}>{\raggedright\arraybackslash}p{0.31\textwidth}>{\raggedright\arraybackslash}p{0.43\textwidth}}
\arrayrulecolor{gray!60}\hline
\rowcolor{grammarheadcolor}$\odot$ & $(\tau_1, \tau_2)$ & $\hat{\odot}(v, v')$ \\
\pyinline{==} & $(\tyObject, \tyObject)$ & $\val{\pyinline{False}}$ if $v$ and $v'$ are both \pyinline{nan}; $\eqOp{\Sigma}(v, v')$ otherwise \\
\pyinline{!=} & $(\tyObject, \tyObject)$ & $\val{\pyinline{True}}$ if $v$ and $v'$ are both \pyinline{nan} or $\eqOp{\Sigma}(v, v') = \val{\pyinline{False}}$; $\val{\pyinline{False}}$ if $\eqOp{\Sigma}(v, v') = \val{\pyinline{True}}$; $\eqOp{\Sigma}(v, v')$ otherwise \\
\pyinline{in} & $(\tyObject, \tyObject)$ & $\inOp{\Sigma}(v, v')$ \\
\pyinline{not in} & $(\tyObject, \tyObject)$ & $\val{\pyinline{True}}$ if $\inOp{\Sigma}(v, v') = \val{\pyinline{False}}$; $\val{\pyinline{False}}$ if $\inOp{\Sigma}(v, v') = \val{\pyinline{True}}$; $\inOp{\Sigma}(v, v')$ otherwise \\
\pyinline{<}\quad\pyinline{<=}\quad\pyinline{>}\quad\pyinline{>=} & $(\tyFloat, \tyFloat)$ & \replaced{$\val{\pyinline{True}}$ if $v < v'$, $v \leq v'$, $v > v'$ or $v \geq v'$ respectively by exact numeric value; $\val{\pyinline{False}}$ otherwise, including when $v$ or $v'$ is \pyinline{nan}}{\todo{Numeric order, as in Python} } \\
\pyinline{<}\quad\pyinline{<=}\quad\pyinline{>}\quad\pyinline{>=} & $(\tyStr, \tyStr)$ & $\val{\pyinline{True}}$ if $v < v'$, $v \leq v'$, $v > v'$ or $v \geq v'$ respectively in the lexicographic order on code points; $\val{\pyinline{False}}$ otherwise \\
\pyinline{+}\quad\pyinline{-}\quad\pyinline{*}\quad\pyinline{//}\quad\pyinline{\%} & $(\tyInt, \tyInt)$ & \replaced{$\val{v + v'}$, $\val{v - v'}$, $\val{v \times v'}$, $\val{\lfloor v / v' \rfloor}$ or $\val{v - v' \times \lfloor v / v' \rfloor}$ respectively; $\aborts{\errZeroDiv}$ for \pyinline{//} or \pyinline{\%} if $v' = 0$}{\todo{Integer arithmetic, as in Python; $\aborts{\errZeroDiv}$ if the divisor is zero} } \\
\pyinline{**} & $(\tyInt, \tyInt)$ & \replaced{$\val{v^{v'}}$ if $v' \geq 0$; $\aborts{\errZeroDiv}$ if $v = 0$ and $v' < 0$; as at $(\tyFloat, \tyFloat)$ otherwise}{\todo{Integer power, as in Python, giving a float if $v'$ is negative} } \\
\pyinline{+}\quad\pyinline{-}\quad\pyinline{*}\quad\pyinline{/}\quad\pyinline{//}\quad\pyinline{\%}\quad\pyinline{**} & $(\tyFloat, \tyFloat)$ & \replaced{An integer operand converted to the nearest float; \pyinline{+}, \pyinline{-}, \pyinline{*} and \pyinline{/} as in IEEE 754 with rounding to nearest, \pyinline{//} and \pyinline{\%} as in Python, and \pyinline{**} with precision left to the implementation; $\aborts{\errZeroDiv}$ for \pyinline{/}, \pyinline{//} or \pyinline{\%} if $v' = 0$, and for \pyinline{**} if $v = 0$ and $v' < 0$; $\aborts{\errOverflow}$ where Python raises it, as for an integer too large to convert}{\todo{Float arithmetic, as in Python, with precision left to the implementation; $\aborts{\errZeroDiv}$ if the divisor is zero, and $\aborts{\errOverflow}$ if Python raises the exception of that name, such as for an integer too large to convert to a float} } \\
\pyinline{+} & $(\tyStr, \tyStr)$ & $\val{v \concat v'}$ \\
\pyinline{+} & $(\tyList{\tyObject}, \tyList{\tyObject})$ & $\val{\exList{\seq{v} \concat \seq{v}'}}$ if $v = \exList{\seq{v}}$ and $v' = \exList{\seq{v}'}$ \\
\pyinline{+} & $(\tyTuple{\tyObject^m}, \tyTuple{\tyObject^n})$ & $\val{\exTuple{\seq{v} \concat \seq{v}'}}$ if $v = \exTuple{v_1, \ldots, v_m}$ and $v' = \exTuple{v'_1, \ldots, v'_n}$ \\
\pyinline{*} & $(\tyStr, \tyInt)$ & $\val{\bigconcat{v^n}}$ where $n = \max(v', 0)$ \\
\pyinline{*} & $(\tyInt, \tyStr)$ & $\val{\bigconcat{v'^n}}$ where $n = \max(v, 0)$ \\
\pyinline{*} & $(\tyList{\tyObject}, \tyInt)$ & $\val{\exList{\bigconcat{\seq{v}^n}}}$ if $v = \exList{\seq{v}}$, where $n = \max(v', 0)$ \\
\pyinline{*} & $(\tyInt, \tyList{\tyObject})$ & $\val{\exList{\bigconcat{\seq{v}^n}}}$ if $v' = \exList{\seq{v}}$, where $n = \max(v, 0)$
\end{tabular}

\vspace{3mm}
\begin{tabular}{>{\raggedright\arraybackslash}p{0.14\textwidth}>{\raggedright\arraybackslash}p{0.31\textwidth}>{\raggedright\arraybackslash}p{0.43\textwidth}}
\arrayrulecolor{gray!60}\hline
\rowcolor{grammarheadcolor}$\ominus$ & $\tau$ & $\hat{\ominus}(v)$ \\
\pyinline{not} & $\tyBool$ & $\val{\pyinline{False}}$ if $v = \pyinline{True}$; $\val{\pyinline{True}}$ otherwise \\
\pyinline{+} & $\tyFloat$ & $\val{v}$ \\
\pyinline{-} & $\tyFloat$ & $\val{-v}$
\end{tabular}
\caption{Run-time overloads of operators}
\label{fig:run-time-overloads}
\end{figure}


The implementation of an operator is selected at run time by the run-time types of its operands, using the
run-time overloads of \figref{run-time-overloads}.
\begin{specdefinition}{Run-time overload}
A \emph{run-time overload} of a binary operator $\odot$ pairs operand bounds $(\tau_1, \tau_2)$ with a definition
of $\hat{\odot}(v, v')$, and applies at operand types
$(\sigma_1, \sigma_2)$ with $\Sigma \vdash \sigma_1 \subty \tau_1$ and $\Sigma \vdash \sigma_2 \subty \tau_2$.
Run-time overloads of a unary operator $\ominus$ pair an operand bound $\tau$ with a definition of $\hat{\ominus}(v)$ in
the same way.
\end{specdefinition}
\noindent The function $\typeOf(v)$ (\S\ref{aux:type-of}) gives the run-time type of a value.
The result $\hat{\odot}(v, v')$ is given by the definition of the most
specific run-time overload of $\odot$ at $(\typeOf(v), \typeOf(v'))$, and is $\aborts{\errExcluded}$ if no run-time
overload of $\odot$ applies (\S\ref{aux:binop}). The result $\hat{\ominus}(v)$ is given in the same way
(\S\ref{aux:unop}). Overload resolution (\secref{well-formedness-operators}) therefore fixes the types of operator
applications but not the run-time overloads. With \pyinline{x: float = 1}, the expression
\pyinline{x + 1} has type \pyinline{float} and evaluates to \pyinline{2} by the run-time overload at
$(\tyInt, \tyInt)$.

\begin{lemma}[Most specific run-time overload]
\label{lem:most-specific-run-time-overload}
For all values $v$ and $v'$, the run-time overloads of $\odot$ at $(\typeOf(v), \typeOf(v'))$ have a least element if
they are not empty, and similarly for $\ominus$ at $\typeOf(v)$, so $\hat{\odot}(v, v')$ and $\hat{\ominus}(v)$ are
defined.
\end{lemma}

\begin{lemma}[Soundness of overload resolution]
\label{lem:overload-resolution-soundness}
If $\vdash v : \sigma_1$, $\vdash v' : \sigma_2$ and $\resolveOp{\Sigma}(\odot, \sigma_1, \sigma_2) = \tau$, then some
run-time overload of $\odot$ applies at $(\typeOf(v), \typeOf(v'))$, and $\vdash v'' : \tau$ if
$\hat{\odot}(v, v') = \val{v''}$; similarly for $\ominus$.
\end{lemma}
\end{AddedBlock}
